\section{A collection that grows one project at a time}
I keep this repository for reusable embedded drivers and for explaining how I build them. Each device has a complete project that can be copied to another directory. The portable library can also be taken alone and connected to another MCU through a small board adapter.

The report is shared documentation, not a shared firmware dependency. There is one English edition and one Persian edition. A new driver becomes a new chapter, with its own hardware, code, acceptance procedure and validation record. The project README remains sufficient for using a project without this report.

\subsection{Current catalog}
\begin{tabular}{p{27mm}p{57mm}p{49mm}}
\toprule
Project & Purpose & Current evidence\\
\midrule
DS18B20 & One-sensor Celsius reading on STM32F107VCT6, PB10 & Builds and portable tests; board acceptance pending\\
STM32F107 foundation & Independent starting project to copy & Consult the template's local status\\
\bottomrule
\end{tabular}

\subsection{Where each part belongs}
\begin{lstlisting}[language={},numbers=none]
ds18b20/
  DS18B20.ioc          # hardware configuration
  Core/               # CubeMX-generated initialization
  Drivers/            # local ST HAL and CMSIS
  Application/        # board adapter and hardware demo
  Libraries/OneWire/  # portable wire protocol
  Libraries/DS18B20/  # portable temperature operation
  MDK-ARM/            # Keil project and startup
  platformio.ini      # GCC build and probe selection
  tests/              # simulated portable C tests
  tools/              # source synchronization/checks
  README.md           # standalone usage instructions
  PROJECT_STATUS.md   # evidence and unfinished work
docs/
  english/            # one complete report edition
  persian/            # matching report edition
  tools/              # report build and PDF verification
\end{lstlisting}

\texttt{Core} initializes the MCU. \texttt{Application} knows the GPIO registers and clock. \texttt{Libraries} contains the reusable C code for talking to the sensor. This separation lets the same temperature code run in a Cortex-M3 simulation or on a different board.

\subsection{How to read and reproduce the work}
For your first result, follow DS18B20 steps 1 through 4 in order. Step 5 teaches the two functions you call in another project; steps 6 and 7 explain the example and pin changes. Steps 8 through 10 explain the implementation after the practical result. You can return to those later. Keep a distinction between a successful compiler run, a passing simulated test and an actual measurement on a board. The final step is always the board acceptance record.

Datasheet PDFs are local working references and are excluded from Git. READMEs and the bibliography preserve official source URLs and revisions. Final report PDFs and editable LaTeX sources are tracked; downloaded fonts, rendered review images and build artifacts are excluded.
